\begin{frame}{이어붙임: 하나의 벡터 안에 두 관점}
  $W_O = I_4$ (연결만 하고 더 이상 섞지 않는 가장 단순한 경우)로
  이어붙이면, 단어 1의 4차원 출력은
  \[ [0.670,\; 0.330,\; 0.330,\; 0.670], \qquad
     \text{단어 2: } [0.330,\; 0.670,\; 0.670,\; 0.330] \]
  읽어보면: 앞 2차원(헤드 A) = ``단어 1의 자기는 $V_1$, 단어 2의
  자기는 $V_2$'' 라는 사실; 뒤 2차원(헤드 B) = ``단어 1의 이웃은
  $V_2$, 단어 2의 이웃은 $V_1$'' 라는 \textbf{다른} 사실.
  단일 헤드의 2차원 출력은 이 두 사실 중 \textbf{하나만} 담을 수
  있었는데, \textbf{멀티헤드는 두 사실을 \emph{같은 벡터 안의
  서로 다른 부분공간}에 나란히 넣었다}.
  \vskip0.3em
  {\footnotesize 실제 학습에서는 ``자기/이웃 헤드''를 누가 맡을지
  \textbf{사람이 설계하지 않는다} — 손실이 줄어드는 방향으로
  역전파가 각 헤드의 $W_i^Q, W_i^K, W_i^V$ 를 \emph{스스로} 다른
  역할로 벌려놓는다.}
  \vskip0.3em
  \kq{그런데 왜 ``합치기(평균/합산)''가 아니라 ``이어붙이기''인가?}
  합하면 $h$개 헤드의 정보가 다시 하나의 벡터로 \textbf{재혼합}되어
  각 헤드의 독립적 관점이 \textbf{희석}된다 — 미러 예에서 숫자로
  확인: 두 헤드 출력을 \emph{더하면}
  \[ [0.670,0.330] + [0.330,0.670] = [1.0, 1.0] \quad
     (\text{단어 2도 } [1.0, 1.0]) \]
  \textbf{두 단어의 출력이 정확히 \emph{같아진다}!} — 미러였던 두
  관점이 합산되는 순간 ``이 단어와 저 단어가 다르다''는 정보가
  \emph{소거}된다. 이어붙임은 각 헤드 출력을 \textbf{서로 다른
  차원(부분공간)}에 두어 소거를 \emph{막고}, $W_O$ 가 ``출력
  차원 $j$는 $i$번째 헤드의 몇번째 성분을 얼마 신뢰할지''를
  결정하게 한다 — ``\textbf{관점의 조합}''을 다음 층(피드 포워드)이
  결정하도록 \textbf{자유도를 남겨둔다}.
  {\footnotesize 실제 학습된 $W_O$ 는 대각선 근방의 ``출력 차원 $j$
  $\leftarrow$ 헤드 $i$의 $k$번째 성분'' 같은 \textbf{희소 패턴}이
  자주 관찰된다.}
\end{frame}
